\chapter{Head Lice}

\section*{Definition}
Pediculosis capitis (head lice infestation) is a highly contagious ectoparasitic infection of the scalp and hair by \textit{Pediculus humanus capitis}, an obligate, blood-feeding insect. It primarily affects children aged 3--11 years and is prevalent worldwide regardless of hygiene.

\section*{What Happens in Your Body}
Adult female lice (2--3 mm) deposit eggs (nits) on hair shafts within 1--2 mm of the scalp, where temperature and humidity are optimal. Eggs hatch in 7--10 days, nymphs mature in 9--12 days, and adults live approximately 30 days feeding on scalp blood 4--5 times daily. Saliva injected during feeding triggers a type IV hypersensitivity reaction, producing pruritus. Transmission requires direct head-to-head contact; fomite transmission (hats, brushes, bedding) is rare.

\section*{Causes}
\begin{itemize}
  \item \textbf{Organism:} \textit{Pediculus humanus capitis} (strictly a human parasite; does not burrow or transmit systemic disease)
  \item \textbf{Transmission:} Direct head-to-head contact; highest risk in school-aged children, households, sleepovers, and camps
  \item \textbf{Risk factors:} Female sex (longer hair, closer contact), age 3--11 years, large household, low socioeconomic status (overcrowding)
  \item \textbf{Resistance:} Permethrin and pyrethrin resistance is widespread in many regions (knockdown resistance mutation kdr)
\end{itemize}

\section*{When to See a Doctor}
See your doctor if:
\begin{itemize}
  \item Persistent scalp itching, especially in children with known exposure
  \item Visible adult lice or nits on hair shafts within 1 cm of the scalp
  \item Secondary bacterial infection of excoriated scalp (impetigo, cervical lymphadenopathy)
  \item No improvement after two complete treatment cycles with a pediculicide
\end{itemize}

\section*{Related Symptoms}
\begin{itemize}
  \item Scalp pruritus, excoriations, cervical lymphadenopathy, pyoderma, impetigo
\end{itemize}
\textbf{Important:} Nits attached more than 1 cm from the scalp are likely non-viable (hatched or dead) and do not require treatment; routine nit removal is recommended for cosmetic and epidemiologic reasons.

\section*{Self-Care / Home Management}
\begin{itemize}
  \item Use a fine-toothed nit comb on wet, conditioned hair every 2--3 days for 2 weeks
  \item Machine-wash bedding, clothing, and hats in hot water (130\textdegree{}F / 54\textdegree{}C) and dry on high heat
  \item Items that cannot be washed: Seal in plastic bags for 2 weeks or freeze at 5\textdegree{}F (--15\textdegree{}C) for 48 hours
\end{itemize}

\section*{How Your Doctor Will Evaluate You}
\begin{itemize}
  \item Clinical diagnosis: Direct visualization of live adult lice; use a fine-toothed comb with conditioner for wet combing
  \item Wood's lamp illumination: Nits fluoresce pearly white (vs. dandruff, hair casts, debris which do not)
  \item Examination of household members and close contacts; all infested individuals should be treated
\end{itemize}

\section*{Treatment Options}
\begin{itemize}
  \item First-line: Permethrin 1\% cream rinse (leave on 10 minutes, repeat day 7--9) if local resistance is low; pyrethrin-based products
  \item Resistance concern: Ivermectin 0.5\% lotion (single 10-minute application), spinosad 0.9\% suspension (single 10-minute application, repeat day 7), or malathion 0.5\% lotion (8--12 hours, flammable)
  \item Oral ivermectin (200 mcg/kg, repeat day 7--9) for refractory cases; not for children <15 kg
  \item Benzyl alcohol 5\% lotion (10-minute application, repeat day 7) for children \ensuremath{\geq}6 months
  \item Wet combing with conditioner alone as a non-chemical alternative (mechanical removal every 3 days for 2 weeks)
\end{itemize}

\clearpage
